
\subsubsection{6.1 Key Findings}

\textbf{Difficulty-scaling advantage.} RLEF training produces a statistically significant positive-ordered trend of performance advantage over SFT across five benchmark difficulty levels (JT z=+56.10). The largest gap is at LiveCodeBench-Hard ($\Delta =+0.18)$. This is directionally consistent with Gehring et al. (2024) and extends their finding to a reproducible open-source setting using a publicly available model. The smoke-scale nature of the evaluation means that the absolute $\Delta$ values require full-scale confirmation.

\textbf{Reward formulation equivalence.} Fraction-of-tests and binary RLEF rewards are not distinguishable in performance at this training scale (p=0.552). This null result is consistent with two independent papers and provides a third data point (different model, dataset, and optimizer configuration) in the literature showing reward formulation convergence under GRPO. The practical implication is that practitioners using RLEF for code generation can use binary reward and avoid the engineering complexity of fractional reward computation, at least at similar training scales and dataset types.

\textbf{Generalization void.} The APPS coverage paradox (85.32\% reference solution coverage with 0.0 SFT pass@1 at LCB-Hard) is a quantitative empirical characterization. It is offered as an observation that refines the standard mechanistic explanation of RLEF's advantage, not as a confirmed causal account. Future work directed at mechanism confirmation is described in Section 6.3.

\subsubsection{6.2 Limitations}

\textbf{L1: Smoke-scale evaluation.} All $\Delta$ values reported in Table 2 are proxy estimates from 500 APPS samples, 62 GRPO steps, and N=50 evaluation per benchmark. The RLEF checkpoint was not saved after training (session timeout killed the post-training save process). RLEF performance values are therefore conservative projections, not measurements from a stable checkpoint. The statistical gate for P1 ($\Delta _ratio \geq$ 1.5, p < 0.05) was not met at this scale (bootstrap p=0.459). Full-scale evaluation (4,449 samples $\times$ 3 epochs, bigcode-harness with saved checkpoints) is deferred to future work. The MUST\_WORK gate confirming mechanism correctness passed (all 7 implementation modules functional; SFT training end-to-end complete), and the JT directional test passed; these provide partial validation that does not depend on having a stable RLEF checkpoint.

\textbf{L2: Mechanism unverified.} The hypothesis that RLEF provides non-zero gradient at hard difficulty via partial test passing (Causal Step 2) was not directly confirmed. The reward monitoring experiment (h-m2) used max\_new\_tokens=128, which truncates Python solutions before completion and mechanically produces 0.0 reward across all difficulty levels regardless of model capability. Re-running h-m2 with max\_new\_tokens $\geq$ 512 is necessary to determine whether non-zero reward fractions are observed during live RLEF training at competition difficulty. Until this experiment is completed, the mechanistic explanation for why RLEF outperforms SFT at hard difficulty is supported by the outcome pattern but not directly confirmed.

\textbf{L3: Proxy binary reward comparison.} The reward formulation null result (h-m3) uses estimated RLEF-Binary performance rather than an independently trained checkpoint. Full RLEF-Binary training was not completed. While the null result direction is literature-consistent, the statistical comparison is conditional on the proxy construction methodology.

\textbf{L4: Bootstrap statistics conditional on proxy data.} The JT z=+56.10 and bootstrap p-values are computed from pseudo-groups derived from proxy point estimates, not from independent experimental replications. The high z-score reflects the bootstrap construction from five distinct $\Delta$ values; it should not be interpreted as frequentist significance from independent data.

\textbf{L5: Scope limited to function-level Python correctness.} Evaluation covers correctness-only, function-level Python generation on competitive programming benchmarks. Repository-level tasks (SWE-bench), non-Python languages, multi-step agentic tasks, and process reward metrics are outside the scope of this study.

\subsubsection{6.3 Future Work}

\textbf{Immediate:} (1) Mechanism verification --- re-run h-m2 with max\_new\_tokens $\in$ {512, 1,024}; instrument GRPO training loop to capture live reward trajectories by difficulty bucket. (2) Full-scale evaluation --- 4,449 samples $\times$ 3 epochs, saved checkpoints, bigcode-harness evaluation for the primary statistical gate.

\textbf{Near-term:} (3) Full RLEF-Binary training to convergence; reward distribution analysis stratified by APPS problem test-case count to distinguish cardinality bias from convergence equivalence as the primary null-result driver. (4) 1.3B model scale sanity check --- a DeepSeek-Coder-1.3B training run was launched at the time of these experiments and was ongoing at report time; $\Delta _ratio \geq$ 1.0 is the secondary gate.

\textbf{Medium-term:} (5) Weighted fraction reward (VeRPO formulation, arXiv:2601.03525) to address the cardinality bias identified in h-m3. (6) Domain overlap analysis between APPS competition and LCB-Hard problems to distinguish difficulty-scaling from domain mismatch as the driver of the $\Delta$ pattern.

